\begin{figure*}[t]
    \centering
    \includegraphics[width=\textwidth]{figures/velox-join-new.pdf}
	\caption{Original execution of a right-inner HashJoin in \proj, which requires converting data from columnar to row-oriented format and back.}
    \label{fig:velox-join}
	\vspace{1em}
\end{figure*}


\begin{figure}[t]
    \raggedright
    \begin{subfigure}{\columnwidth}
        \raggedright
        \includegraphics[width=\columnwidth]{figures/rowcontainer_layout.pdf}
        \caption{Original row-oriented layout in \proj, where key and payload columns are co-located within each row.}
        \label{fig:rowcontainer-layout}
    \end{subfigure}

    \vspace{0.6em}

    \begin{subfigure}{\columnwidth}
        \raggedright
        \includegraphics[width=\columnwidth]{figures/hybridcontainer_layout.pdf}
        \caption{Hybrid layout in \proj, where key columns are stored in a row-oriented container and payload columns remain in separate columnar storage referenced via a compact \emph{rowRef}.}
        \label{fig:hybridcontainer-layout}
    \end{subfigure}

    \caption{Comparison of row layouts in \proj.}
    \label{fig:rowcontainer-layouts}
\end{figure}







\section{Hybrid Execution Model}
\label{sec:layout}
\subsection{Existing Design and Limitations}

In existing vectorized execution engines, operators frequently switch between columnar and row-oriented representations during query execution, depending on which layout best supports a given computation.
For example, in \proj~\cite{bolt} and Velox~\cite{velox}, a HashJoin typically converts its columnar input into a row-oriented format before building the hash table, and reconstructs columnar output after probing.
Figure~\ref{fig:velox-join} illustrates this process for a right inner HashJoin, where the build side contains both join keys and payload columns, while the probe side provides only join keys.

On the build side, input columnar vectors are packed into a \emph{RowContainer} that stores each tuple contiguously in memory.
A hash table is built on the join keys, with each entry pointing to a corresponding row in the \emph{RowContainer}.
During probing, join keys from the probe-side input are used to probe the hash table.
For each match, the stored row pointer is dereferenced, and the projected payload fields are extracted and written into output columnar vectors.
Algorithm~\ref{algo:velox-extract} illustrates this materialization process.

A similar design is used by Sort.
Input columnar vectors are first converted into row-oriented tuples stored in a \emph{RowContainer}, and sorting is performed by comparing key fields via row pointers.
After sorting completes, the engine reconstructs the sorted row pointers in columnar form to conform to the vectorized execution interface.

Figure~\ref{fig:rowcontainer-layout} shows the internal layout of the \emph{RowContainer}.
Each row stores all serialized key fields contiguously, followed by their null bitmap.
Payload (non-key) fields are stored next, also accompanied by a null bitmap.
For HashJoin, rows additionally contain a next-chain pointer that links tuples with identical join keys, enabling efficient traversal during probing.

While this design improves locality for key access during hashing, probing, and comparison, it has several important limitations.
First, operators eagerly convert entire tuples into a row-oriented format, even though key-centric operators typically access only a small subset of columns.
This forces unnecessary copying and movement of payload data that is not immediately needed.
Second, storing keys and payloads together in wide rows reduces cache efficiency for key-centric access, as cache lines fetched for key processing also bring in unused payload fields.
Finally, the tightly coupled row-oriented layout prevents operators from optimizing the representation of keys independently from payloads, limiting opportunities for fine-grained layout specialization.
These limitations motivate the hybrid execution model introduced in the next section.

\begin{algorithm}[t]
\SetKwInOut{Input}{Input}
\Input{
  $C$: set of projected column indices;\\
  $R$: list of matched row pointers;\\
  $O$: output RowVector
}
\ForEach{$c \in C$}{
  \For{$i \gets 0$ \KwTo $|R|-1$}{
    $r \gets \textsc{RowContainer.getRow}(P[i])$\;
    $O.\textsc{childAt}(c)[i] \gets r.\textsc{valueAt}(c)$\;
  }
}
\caption{Original columnar materialization from a row-oriented layout in \proj.}
\label{algo:velox-extract}
\end{algorithm}


\begin{figure*}[t]
    \centering
    \includegraphics[width=\textwidth]{figures/hybrid-join-new.pdf}
	\caption{Hybrid execution of a right-inner HashJoin in \proj, where only key columns are converted to a row-oriented layout.}
    \label{fig:hybrid-join}
	\vspace{1em}
\end{figure*}





\subsection{Our Approach}

Our key observation is that many operators require row-wise access only for their \emph{key columns}, while \emph{payload columns} are accessed primarily to produce final results.
For example, in joins and sorts, key columns are repeatedly accessed together during processing (e.g., for hashing or comparison), whereas payload columns are accessed only during output materialization.
Despite this asymmetry, existing designs apply layout conversions uniformly to all columns, forcing both keys and payloads into the same row-oriented representation.
This leads to substantial CPU and memory overhead, particularly for wide tables and complex query plans where the same data may undergo repeated conversions across operators.

To address this inefficiency, we introduce a \textbf{hybrid execution model} that processes different parts of a tuple using the layout best suited to their access patterns.
Specifically, key columns are stored and accessed in a row-oriented structure to enable efficient key-centric processing, while payload columns remain in their original columnar representation to avoid unnecessary memory copying and layout conversions.

In \proj, this design is implemented as a \emph{HybridContainer}, which combines a row-oriented \emph{RowContainer} for key storage with columnar \emph{RowVector}s for payload storage.
Figure~\ref{fig:hybrid-join} illustrates the workflow of the hybrid execution model for a HashJoin.
On the data-generating side (HashBuild), key columns from the input vectors are packed into the \emph{RowContainer} to maximize cache locality during hash table construction, while payload columns are retained in columnar \emph{RowVector}s extracted directly from the input.
Each row stored in the \emph{RowContainer} additionally records a compact \emph{rowRef}, which serves as a logical reference to the corresponding payload tuple in columnar storage.

During probing, join keys from the probe-side input are used to probe the hash table.
For each match, the stored \emph{rowRef} identifies the payload tuple associated with the matched key, allowing the engine to retrieve payload columns directly from columnar storage without packing them into rows.
The output is then produced in columnar form by combining key columns retrieved from the \emph{RowContainer} with payload columns accessed via \emph{rowRef}.
A similar execution pattern is used by Sort: after sorting the list of row references based on key comparisons, key columns are retrieved from the \emph{RowContainer}, while payload columns are reconstructed using the corresponding \emph{rowRef}s.

Figure~\ref{fig:hybridcontainer-layout} illustrates the internal memory layout of a row stored in the \proj \emph{HybridContainer}.
Each row still begins with serialized key values and their associated null bitmap, followed by a 64-bit \emph{rowRef}, and an optional chaining pointer used by hash-based operators. Compared to the original row-oriented layout, this compact row representation reduces the cache-line footprint of each row and improves cache-line utilization during key-centric processing, allowing more rows to fit in cache and improving access efficiency for hashing, comparison, and traversal.

The \emph{rowRef} is a compact, immutable identifier that records where the payload of a tuple is stored in columnar form. 
In vectorized engines such as \proj and Velox, query execution is organized into pipelines that may be executed by multiple parallel \emph{drivers}~\cite{velox-pipelines}.
Each driver processes a disjoint subset of the input and maintains its own local operator state.
For multi-driver operators such as HashJoin, rows produced by different drivers cannot be uniquely identified by a local row index alone.
Accordingly, \emph{rowRef} encodes both the identifier of the data-generating operator instance (e.g., the HashBuild operator in a HashJoin) and the tuple’s row index within that operator’s input.
In our implementation, \emph{rowRef} is a 64-bit value that dedicates the high-order bits (by default, 8~bits) to the driver identifier and the remaining bits to the row index, enabling a compact and globally unique reference across all parallel drivers.

This representation uniquely identifies the origin of each tuple and enables efficient, direct access to payload columns during output materialization, as illustrated in Algorithm~\ref{algo:hybrid-materialize}.
Overall, the hybrid execution model avoids eagerly converting payload columns into row-oriented formats, substantially reducing layout conversion and memory copying overheads while preserving cache locality for key-intensive operations such as HashJoin and Sort.
As shown in Figure~\ref{fig:join-breakdown}, without affecting hash table build performance, the hybrid execution makes the input conversion phase $5.8\times$ faster and improves overall HashJoin performance by up to $1.7\times$.



\begin{algorithm}[t]
\SetKwInOut{Input}{Input}
\Input{
  $C$: set of projected column indices;\\
  $R$: list of matched row references;\\
  $O$: output RowVector;\\
  $B$: array of payload RowVectors (columnar batches)
}
\ForEach{$c \in C$}{
  \uIf{$\textsc{isKey}(c)$}{
    \For{$i \gets 0$ \KwTo $|R|-1$}{
      $r \gets \textsc{RowContainer.getRow}(R[i])$\;
      $O[c][i] \gets r.\textsc{valueAt}(c)$\;
    }
  }
  \Else{
    \For{$i \gets 0$ \KwTo $|R|-1$}{
      $(did, rid) \gets \textsc{decodeRef}(R[i])$\;
      $O.\textsc{childAt}(c)[i] \gets B[did][c][rid]$\;
    }
  }
}
\caption{Hybrid columnar materialization in \proj: key columns are extracted from the \textsc{RowContainer}, while payload columns are read directly from columnar storage.}
\label{algo:hybrid-materialize}
\end{algorithm}




% \noindent\textbf{Analysis.}The proposed hybrid in-memory execution model offers two main advantages.
% First, it eliminates the need for memory copying and full column-to-row conversion for key‑intensive operators, reducing both CPU overhead and memory allocation costs.
% Second, by separating key columns from payload columns, it improves CPU cache utilization: fetching keys no longer pulls unnecessary payload values into the cache, reducing cache misses during key-intensive operations such as HashBuild and HashProbe.

% The design also introduces trade-offs.
% Each row stores an additional \texttt{rowId}, which incurs extra memory overhead.
% This design is a well-known technique in columnar database systems, used for late materialization in systems such as MonetDB~\cite{boncz2005monetdb} and HyPer~\cite{kemper2011hyper} to reconstruct tuples from separate columnar storage.
% In our case, the \texttt{rowId} serves the same purpose entirely in memory, enabling payload retrieval without storing it in row format.
% Output materialization thus requires a rowId-based lookup into columnar batches, adding a small CPU cost; for example, Figure~\ref{fig:join-breakdown} shows that the hybrid design incurs about 15\% additional CPU time in HashProbe.
% In practice, we observe that real‑world queries often have wide payload columns with a limited key columns (for example, a Sort on 4 key columns over a table with around 50 projected columns), making the overhead of storing \texttt{rowId} negligible compared to the performance gains.


\subsection{Limitations and Optimizations}

The hybrid execution model reduces \textbf{intra-operator} layout conversion overheads by separating key and payload processing.
By storing keys in a row-oriented layout and retaining payloads in columnar form, it avoids eagerly materializing wide tuples into rows while preserving cache-friendly, key-centric access for operations such as hashing, comparison, and partitioning.

These benefits have been discussed in detail in the preceding subsections; here we focus on the limitations of the design and the optimizations used to mitigate them.

\medskip
\noindent\textbf{Limitations.}
The primary cost introduced by the hybrid design arises during output materialization.
Unlike the original implementation, which reconstructs output columns from a row-oriented container, the hybrid design retrieves payload values via an indirect reference (\emph{rowRef}) into columnar storage.

In the original materialization procedure (Algorithm~\ref{algo:velox-extract}), rows are stored contiguously in the \emph{RowContainer}.
Although materialization proceeds column by column, extracting multiple columns for the same set of rows benefits from good cache locality: cache lines fetched while accessing one column of a row are often reused when accessing subsequent columns of the same row.
As a result, the original design exhibits relatively cache-friendly behavior on the read path during materialization.

In contrast, the hybrid materialization procedure (Algorithm~\ref{algo:hybrid-materialize}) accesses payload columns through \emph{rowRef}, which may refer to tuples originating from different drivers and columnar batches.
These accesses are inherently less cache-friendly, as payload rows are scattered across columnar vectors, leading to additional cache misses and TLB pressure.
In our evaluation, this manifests as a more than 15\% increase in CPU time during tuple reconstruction phases—such as the HashProbe stage of HashJoin and the getOutput stage of Sort—when no optimizations are applied.

In addition, the hybrid design is most effective for operators with wide payloads.
When the projected payload is narrow (e.g., a single \texttt{BIGINT} column), storing payload values directly alongside keys in a row-oriented layout is already efficient.
In such cases, replacing an 8-byte payload value with an 8-byte \emph{rowRef} yields little benefit for key-centric processing, while introducing additional indirection during output materialization.
As a result, the hybrid layout can underperform the original design for queries with very small payloads.

For queries with moderate to wide payloads, however, the hybrid design avoids eagerly materializing large payloads into rows prior to operator execution.
In this workload setting, the cost of indirect payload access during output reconstruction is outweighed by the savings from eliminating repeated column-to-row layout conversions, resulting in net performance improvements.

The memory overhead introduced by the hybrid design is modest in the target workload setting.
Each row stores an additional 64-bit \emph{rowRef}, corresponding to 8~bytes per row.
For wide analytical workloads (e.g., more than 20 payload columns and average tuple widths exceeding 200~bytes), this overhead accounts for less than 5\% of the total in-memory footprint.



\medskip
\noindent\textbf{Optimizations.}
To account for the workload-dependent trade-offs described above, \proj selectively enables the hybrid execution model only when it is expected to be beneficial.
Specifically, we introduce a configurable threshold on the total payload width: the hybrid execution model is enabled only when the projected payload size exceeds a predefined number of bytes.
In our implementation, the default threshold is set to 24~bytes (equivalent to three \texttt{BIGINT} values).

Beyond this adaptive selection, several implementation-level optimizations further mitigate the cost of indirect payload access.
Our implementation in \proj incorporates the following techniques:

\begin{itemize}[leftmargin=*]
  \item \emph{Payload batch coalescing}: After all input has been consumed, columnar payload batches are merged into a smaller number of contiguous batches to reduce TLB misses during extraction.
  \item \emph{Driver-aware materialization}: For multi-driver operators such as HashJoin, matched rows are grouped by driver identifier prior to extraction, improving cache locality when accessing columnar payloads.
  \item \emph{Software prefetching}: Payload accesses are prefetched during materialization using a fixed prefetch distance and loop unrolling (e.g., four-way unrolling) to overlap memory latency with computation.
\end{itemize}

We evaluate the effectiveness of these optimizations in Section~\ref{sec:experiments} by comparing hybrid execution with and without these optimizations enabled.

\medskip
\noindent\textbf{Applicability to other operators and systems.}
Although implemented in \proj for HashJoin and Sort, the hybrid execution model naturally extends to other key-centric operators such as HashAggregation and Window.
More broadly, the design is applicable to modern vectorized engines that rely on row-oriented representations for key processing, including systems such as Velox, Presto, DuckDB, and ClickHouse.
The hybrid abstraction thus represents a general architectural pattern for reducing intra-operator layout conversion overheads in analytical query processing.





% \paragraph{Advantages.}
% This design offers several benefits:  
% (i) it reduces the memory footprint of the \texttt{RowContainer};  
% (ii) it eliminates unnecessary columnar-to-row and row-to-column conversions for payload data;  
% (iii) it improves CPU efficiency by maintaining cache-local key access; and  
% (iv) it naturally supports our payload-first adaptive spilling strategy described in Section~\ref{sec:adaptive-spill}.




% Velox adopts a purely columnar in-memory layout for intermediate data, representing each column as a contiguous \texttt{Vector} to enable efficient vectorized execution. While this layout excels for scan-heavy operators, certain operators—most notably HashJoin and Sort—require accessing multiple columns from the same row or maintaining per-row state in memory. To support these operators, Velox performs a \emph{column-to-row conversion} by packing selected columns into a \texttt{RowContainer}, a row-oriented storage abstraction optimized for random access to complete rows.

% This conversion, while enabling row-oriented processing, incurs non-trivial CPU and memory overhead. For example, in a representative TPC-H query with a large build-side join table (\texttt{HashBuild}), profiling shows that \texttt{Column-to-Row} conversion accounts for \textbf{XX\%} of the total build time, while the actual hash table insertion accounts for \textbf{YY\%}. Figure~\ref{fig:hashbuild-breakdown} presents a detailed breakdown, illustrating that layout conversion alone can consume a significant fraction of the operator’s runtime.

